% =====================================================================
%  FICHE 11 — THEME 3 — Exercices FACILE — Équilibrer une équation redox
% =====================================================================
\documentclass[11pt,a4paper]{article}
\usepackage[utf8]{inputenc}
\usepackage[T1]{fontenc}
\usepackage{lmodern}
\usepackage[french]{babel}
\usepackage{amsmath,amssymb}
\usepackage[version=4]{mhchem}
\usepackage{siunitx}
\sisetup{output-decimal-marker={,},per-mode=symbol}
\usepackage{xcolor}
\usepackage{array}
\usepackage{booktabs}
\usepackage{enumitem}
\usepackage[most]{tcolorbox}
\usepackage[a4paper,margin=2.1cm,headheight=26pt,headsep=10pt]{geometry}
\usepackage{fancyhdr}
\usepackage[hidelinks]{hyperref}

\definecolor{primary}{RGB}{146,95,160}
\definecolor{primarydark}{RGB}{92,52,104}
\newcommand{\NO}{\ensuremath{\mathrm{N.O.}}}

\newcounter{exo}
\newtcolorbox{exo}[1][]{enhanced,breakable,boxrule=0.9pt,arc=2pt,colback=primary!4,
  colframe=primary,fonttitle=\bfseries,coltitle=white,left=3mm,right=3mm,top=2.2mm,bottom=2.2mm,
  title={Exercice~\refstepcounter{exo}\theexo\ifx\relax#1\relax\relax\else\ \textmd{--- \itshape #1}\fi}}

\pagestyle{fancy}\fancyhf{}
\fancyhead[L]{\small\textcolor{primarydark}{\textbf{Fiche 11 · Thème 3} — Équilibrer une équation redox}}
\fancyhead[R]{\small\textcolor{primarydark}{\textsc{Facile}}}
\fancyfoot[C]{\small\thepage}
\renewcommand{\headrulewidth}{0.8pt}
\renewcommand{\headrule}{\hbox to\headwidth{\color{primary}\leaders\hrule height \headrulewidth\hfill}}

\begin{document}
\begin{center}
{\LARGE\bfseries\color{primarydark}Thème 3 — Niveau Facile}\\[2pt]
{\small\itshape Équilibrer une seule demi-équation en milieu acide. Procède O, puis H, puis électrons.}\\[-4pt]
{\color{primary}\rule{\textwidth}{1pt}}
\end{center}

\begin{exo}[la demi-équation du permanganate — milieu acide]
Équilibre en milieu acide la demi-équation de réduction du couple \ce{MnO4^-}/\ce{Mn^{2+}} :
\[
\ce{MnO4^- -> Mn^{2+}}.
\]
Détaille les trois étapes : équilibrage de l'oxygène (\ce{H2O}), de l'hydrogène (\ce{H+}), puis de
la charge (\ce{e^-}).
\end{exo}

\begin{exo}[le dichromate — attention au facteur 2]
Équilibre en milieu acide la demi-équation de réduction du couple \ce{Cr2O7^{2-}}/\ce{Cr^{3+}} :
\[
\ce{Cr2O7^{2-} -> Cr^{3+}}.
\]
Commence par équilibrer l'élément chrome : combien d'atomes de \ce{Cr^{3+}} obtiens-tu ? Combien
d'électrons au total ?
\end{exo}

\begin{exo}[le nitrate en monoxyde d'azote]
Équilibre en milieu acide la demi-équation de réduction du couple \ce{NO3^-}/\ce{NO} :
\[
\ce{NO3^- -> NO}.
\]
\end{exo}

\begin{exo}[deux demi-équations sans oxygène]
Équilibre les deux demi-équations suivantes (aucun atome d'oxygène : seuls les électrons sont à
placer) :
\[
\ce{Fe^{2+} -> Fe^{3+}} \qquad\text{et}\qquad \ce{Sn^{2+} -> Sn^{4+}}.
\]
Précise dans chaque cas s'il s'agit d'une oxydation ou d'une réduction.
\end{exo}

\begin{exo}[le contrôle des charges]
On propose la demi-équation équilibrée :
\[
\ce{MnO4^- + 8 H+ + 5 e^- -> Mn^{2+} + 4 H2O}.
\]
Vérifie qu'elle est correcte en calculant la \textbf{somme des charges} de chaque côté. Vérifie aussi
la conservation des atomes de Mn, O et H.
\end{exo}

\end{document}
